\chapter{Il plug troncato sotto il ciclo}
\label{ch:plugcycle}

Il Capitolo~\ref{ch:plug} sa marciare un plug a una condizione di
camera. Un motore a detonazione rotante non offre mai una sola
condizione: la sua onda di detonazione spazza l'anello e una sezione
fissa vede la pressione di camera salire di colpo e svuotarsi,
ancora e ancora. Questo capitolo taglia corta la spina --- i plug
reali sono sempre troncati, perché la coda lunga e sottile pesa
moltissimo e produce quasi nessuna spinta --- e si chiede dove
tagliare quando la pressione di camera oscilla in quel modo. La
risposta non risulta essere ``progettare per la condizione media'', e
la differenza è misurabile.

\section{La domanda}

Due tipi di ugello si comportano in modo opposto su un ciclo. La
campana fissa \emph{media}: il suo ottimo di ciclo è il progetto per
la singola condizione di pressione media, che il motore ha
certificato nel Capitolo~\ref{ch:cycle}. Il plug completo invece
\emph{massimizza}: l'ottimo di ogni fase contiene ogni plug più
lungo, quindi il plug ideale è dimensionato sulla fase più forte. Il
plug \emph{troncato} è il punto in cui entrambi i meccanismi
falliscono insieme --- ``il primo problema di forma genuinamente
mediato''~\cite{theory}. Il motore sa misurare tutte e tre le
affermazioni: l'ordinamento, la sua rottura e la condizione di
ottimalità mediata che la sostituisce?

\section{Come funziona, per esteso}

\subsection{Che cosa afferma la teoria, e che cosa assume}

Prima due termini. Una \emph{fase di ciclo} $\xi \in [0,1]$ indica
dove ci si trova nella storia di camera che si ripete --- $\xi$
vicino al passaggio dell'onda è la parte ad alta pressione, $\xi$
vicino alla fine dello svuotamento è quella a bassa pressione
(Capitolo~\ref{ch:cycle}); la media di ciclo $\mu$ è la media di
tipo temporale su quelle fasi. E un plug è \emph{troncato} alla
lunghezza $l$ quando la spina è tagliata lì e chiusa con una base
piatta: l'intera questione di progetto è quale debba essere $l$.

Il risultato del corpus sui plug (la sua T-T4) è dimostrato sotto
una chiusura di \emph{adattamento ideale} --- un'assunzione fatta per
mantenere trattabile la dimostrazione: che a valle del punto di
piena espansione di ciascuna fase la pressione di parete si agganci
esattamente all'ambiente $p_a$, e a monte di esso scali con la
pressione di camera. Sotto quell'assunzione la spinta di una data
fase, $F(l;\xi)$, non decresce mai al prolungarsi del plug e diventa
esattamente costante oltre una lunghezza dipendente dalla fase
$l(\xi)$, con $l(\xi)$ maggiore per le fasi più forti. L'insieme
delle lunghezze ottimali di ciascuna fase è allora una semiretta
$[l(\xi), \infty)$, e queste semirette sono \emph{annidate}: quella
della fase più forte sta dentro tutte le altre. È l'annidamento a
far collassare il problema --- una sola lunghezza è ottimale per
tutti insieme, quindi la migliore media eguaglia la media dei
migliori,
$\max_\Sigma \int F\,\mathrm{d}\mu = \int \max F\,\mathrm{d}\mu$,
raggiunta dal plug non troncato dimensionato sulla fase di picco.

L'avvertenza di finezza allegata a quel teorema dice che cosa
distrugge l'annidamento: un limite sulla lunghezza, un modello di
pressione di base o un adattamento non ideale. Perso l'annidamento,
l'uguaglianza diventa una disuguaglianza stretta,
$\max\int < \int\max$, e l'ottimo soddisfa invece una
\emph{condizione di spigolo mediata} sul piano di troncamento --- i
voti di pressione di parete per fase devono lì bilanciarsi nella
media di ciclo. Ciò che il motore aggiunge è precisamente la
quantità che la chiusura elimina per ipotesi: la pressione di parete
realmente marciata non si aggancia. A valle dell'adattamento va
sotto e sopra $p_a$, e l'ampiezza di quello scostamento \emph{è} la
rottura dell'adattamento, misurabile fase per fase.

\subsection{Mediare il problema non è mediare la condizione}
\label{sec:avgnotavg}

La media di ciclo $\mu$ è una media di tipo temporale sulle fasi, ed
è tentante concluderne che la ``condizione di spigolo mediata su
$\mu$'' sia solo il progetto alla pressione media travestito. Non lo
è, e la distinzione è l'intero contenuto della dicotomia. Il
\emph{progetto alla pressione media} significa: costruire il singolo
flusso stazionario dello \emph{stato} mediato --- la fase la cui
pressione di camera è $\bar P_0$ --- e ottimizzare quell'unico
flusso. Per la campana fissa questo è legittimo, perché il flusso e
la spinta della campana sono \emph{lineari} nella pressione di fase:
la legge di scala (Lemma~A) dice che l'intero campo scala con $P_0$
su una geometria fissa, con l'ambiente che entra solo in modo
additivo attraverso il gauge di pressione. Per un problema lineare la
media delle soluzioni è la soluzione della media --- mediare-poi-
risolvere equivale a risolvere-poi-mediare. È esattamente per questo
che la campana collassa.

La \emph{condizione di spigolo mediata} invece risolve ogni fase come
il \emph{proprio} flusso non lineare, lascia che ciascuna esprima il
proprio voto $p_w(l^{*};\xi) - p_a$ sul piano di troncamento ---
positivo significa che quella fase sta ancora spingendo lì e vuole
più plug, negativo che viene trascinata e ne vuole meno --- e media
la \emph{condizione di ottimalità}:
$\int (p_w(l^{*};\xi) - p_a)\,\mathrm{d}\mu = 0$. Le due coincidono
solo se l'applicazione da fase a voto è affine --- e il plug rompe
quell'affinità attraverso il suo bordo libero: l'ambiente è
\emph{condiviso e fisso}, non scala con la fase, quindi la pressione
di parete di ogni fase satura verso lo \emph{stesso} $p_a$ nel
\emph{proprio} punto di adattamento. La media dei voti non è il voto
alla media. La suite qui sotto misura la differenza direttamente: il
progetto alla pressione media dà $l = 1.280$~m, la condizione mediata
colloca l'ottimo a $l^{*} = 2.103$~m --- tutt'altro luogo. In una
riga: il risultato della campana media il \emph{problema}; quello del
plug media la \emph{condizione}, e la media commuta solo con i
problemi lineari.

\subsection{Una marcia per fase copre ogni troncamento}

In flusso supersonico le perturbazioni non possono risalire la
corrente: ciò che accade in un punto è fissato interamente da una
regione limitata alle sue spalle, il suo \emph{dominio di
dipendenza}. Su una spina \emph{fissa}, troncare a $l$ rimuove
parete solo a valle di $l$, quindi nulla a $x < l$ può sentire il
taglio. Una sola marcia a spina completa per fase produce dunque in
forma chiusa l'intera famiglia dei troncamenti,
\begin{equation}
J(l;\xi) \;=\; F_{\rm in}(\xi)
\;+\; \int_{x_0}^{l} \left(p_w(x;\xi) - p_a\right)
\,2\pi y_w\,(-\mathrm{d}y_w),
\label{eq:trunc}
\end{equation}
dove ogni pressione è misurata rispetto all'ambiente, così il bordo
libero contribuisce zero (la pressione su di esso \emph{è} $p_a$) e
anche un piano di base a $p_b = p_a$ contribuisce zero. Derivando in
$l$ si ottiene la condizione marginale $\mathrm{d}J/\mathrm{d}l = 0$,
cioè $p_w(l) = p_a$ --- la condizione di spigolo del plug sul piano
di troncamento: estendere il plug esattamente fin dove la sua parete
sta ancora spingendo più forte dell'atmosfera. La sua media di ciclo
sulle fasi è la condizione di spigolo mediata di cui sopra. Lo studio
dei troncamenti a cinque fasi costa perciò cinque marce, e ogni
verifica qui sotto legge le curve di pressione di parete
memorizzate.

\subsection{Il mondo}

Assialsimmetrico, tabelle NASA di gas reale, un unico pezzo di
hardware fisso marciato da tutte le fasi. Labbro del cowl in
$(0, 2)$; il flusso anulare arriva a $M_i = 1.2$ diretto verso
l'asse ($\theta_i = -25^{\circ}$); il ventaglio di labbro ---
l'espansione di spigolo del Capitolo~\ref{ch:plug}, costruita per
ogni fase e coerente con le tabelle --- lo espande fino all'ambiente
condiviso. L'ambiente è posto in modo che la rotazione del ventaglio
sia circa $|\theta_i|$, cioè che il flusso esca all'incirca in
direzione assiale ($\theta_E \approx 0$): la classica condizione di
progetto del plug, che rende la spina --- una linea di corrente del
campo del ventaglio della fase media --- genuinamente
\emph{discendente} verso l'asse. La linea iniziale pone lo stesso
profilo $(M,\theta)(y)$ per ogni fase con la termodinamica propria
di ciascuna fase: dati esatti da legge di scala per il ciclo a sola
pressione, e la posizione naturale ``stesso profilo di Mach
all'iniettore'' per il ciclo simulato. Ogni fase gira attraverso la
marcia del plug certificata del Capitolo~\ref{ch:plug} (parete
assegnata, ricerca multi-colonna del piede, contabilità delle
righe); il ciclo stesso è la costruzione certificata del
Capitolo~\ref{ch:cycle} (cinque fasi log-uniformi, rapporto di
pressione $PR = 10$, rapporto di temperatura simulato
$TR = 1.306$).

\section{Che cosa è stato costruito}

\carrier{validation/a1\_plug\_cycle.py}: il costruttore del mondo
(ventaglio per fase coerente con le tabelle, linea di corrente della
spina sul campo medio), i kit di marcia per fase con cache su disco,
la famiglia di troncamenti in forma chiusa~\eqref{eq:trunc} e due
stadi --- \code{t3} (ciclo a sola pressione: la legge di scala a
livello di parete e l'anti-collasso a livello di troncamento) e
\code{mock} (profili di $P$ e $T$: l'ordinamento degli ottimi per
fase, la rottura misurata dell'adattamento, i voti di spigolo
mediati, il divario stretto $\int\max > \max\int$, la penalità del
progetto di picco) --- più un rigettatore ad ambiente corrotto.

\section{Risultati di riferimento}

\textbf{Ciclo a sola pressione (l'anti-collasso del plug): 5/5
PASS.} Tutte le marce di fase certificate da Newton. La legge di
scala a livello di parete compare dal vivo esattamente dove deve e
soltanto lì: sulla parte di parete ancora schermata dall'ambiente ---
a monte del primo arrivo del bordo, $x < 0.81$, trovato tracciando
all'indietro la $C^{-}$ del bordo fino alla parete --- le pressioni
di parete scalate $p_w/P_0$ collassano su un'unica curva a meno di
$4.7\times10^{-15}$, lo zero di macchina. Oltre quella stazione il
collasso si rompe, e quella rottura \emph{è} il meccanismo di
adattamento, nettamente localizzato. I troncamenti misurano poi la
dicotomia: la lunghezza ottimale di ciclo $l^{*} = 2.103$ sta molto
\emph{sopra} il progetto alla pressione media ($l = 1.280$) e appena
sotto il progetto di picco ($l = 2.336$) --- il plug \emph{non}
media, si sposta verso il picco, trattenuto al di qua di esso solo
dalla reale rottura dell'adattamento. Gli ottimi per fase
$l(\xi) = 2.336, 2.061, 1.125, 1.047, 1.006$ crescono con la
pressione di fase: l'\emph{ordine} di annidamento del teorema del
plug, dal vivo.

\textbf{RDE simulato (profili di $P$ e $T$), la media genuina: 9/9
PASS.} Lo stesso ordinamento ricompare
($l(\xi) = 2.336 \ldots 1.012$). L'adattamento non ideale è
\emph{misurato}: ogni fase interna perde spinta oltre il proprio
$l(\xi)$ (perdite relative da $1.6\times10^{-4}$ a
$2.0\times10^{-3}$) --- l'aggancio ideale non vale nel flusso reale,
e quella perdita è esattamente la quantità che rompe l'annidamento
ideale. Di conseguenza la disuguaglianza stretta dell'avvertenza di
finezza è risolta:
$\int\max - \max\int = 1.07\times10^{4}$ N ($2.3\times10^{-4}$ di
$J$), oltre la sua banda derivata; la penalità per progettare al
picco è $2.2\times10^{-4}$ di $J$, anch'essa risolta. All'ottimo di
ciclo trovato alla cieca la condizione di spigolo mediata vale: la
media dei voti è $+1.8\times10^{4}$ Pa contro una banda di
equilibrio di $3.1\times10^{5}$, mentre i singoli voti raggiungono
$\pm1.4\times10^{6}$ --- e l'ottimo è un compromesso strettamente
interno (gli ottimi per fase si distribuiscono su 1.32 in $l$;
nessuna fase estrema possiede $l^{*}$). Un ambiente corrotto sposta
la media dei voti $2.3\times$ fuori dalla banda --- rigettato. La
strada per arrivarci --- due lezioni di posizione del problema e due
verifiche che la teoria stessa ha rigettato --- è
l'Appendice~\ref{app:plugcycleposing}.

\textbf{Una lezione misurata che è fisica nuova, non un difetto.} La
struttura puntuale dei voti della campana --- monotona in fase, con
esattamente un cambio di segno (Capitolo~\ref{ch:cycle}) --- \emph{non}
si trasferisce al plug. A valle dell'adattamento la pressione di
parete \emph{oscilla} attorno a $p_a$, perché il ventaglio di labbro
si riflette sul bordo libero e ricade sulla parete. I voti puntuali
campionano dunque un'oscillazione, e la fase di passaggio dell'onda
può votare negativo a una lunghezza inferiore al proprio ottimo.
Sopravvivono solo il bilancio mediato e la struttura di compromesso
--- che è esattamente, e soltanto, ciò che l'avvertenza di finezza
afferma.

\begin{figure}[htbp]
\centering
\includegraphics[width=\textwidth]{fig_plug_cycle.pdf}
\caption{Il plug troncato sotto il ciclo, misurato. (a)~Pressione di
parete lungo la spina fissa per ciascuna fase: $x$ [m] contro $p_w$
[bar], con l'ambiente condiviso $p_a$ disegnato come linea
orizzontale. Si guardi a valle della regione schermata: le curve non
si agganciano a $p_a$, vi \emph{oscillano} attorno --- il ventaglio
di labbro che si riflette sul bordo libero e torna sulla parete.
Quell'oscillazione è l'adattamento reale, non ideale, che la chiusura
della teoria elimina per ipotesi. (b)~Le famiglie di troncamento
dell'Equazione~\eqref{eq:trunc}: lunghezza di troncamento $l$ [m]
contro la deviazione di spinta di ciascuna curva dal proprio massimo,
$J(l;\xi) - \max_l J(l;\xi)$, in kN, così che zero significhi
``all'ottimo proprio di questa fase''. I massimi per fase ($\nabla$)
si distribuiscono da $l = 1.01$ a $2.34$~m nell'ordine della
pressione di fase, e ogni curva \emph{cala} oltre il proprio massimo
--- la rottura misurata dell'annidamento ideale, che le avrebbe
tenute piatte. La media di ciclo $J_{\rm cycle}(l)$ (nero) ha il
picco a un compromesso $l^{*}$ che nessuna singola fase possiede.
(c)~I voti a quell'ottimo: fase di ciclo $\xi$ [--] contro
$p_w(l^{*};\xi) - p_a$ [bar]. Sono scompaginati nell'ordine dei segni
dall'oscillazione di adattamento --- la struttura monotona dei voti
della campana non si trasferisce --- eppure la loro media di ciclo
(tratteggiata) si bilancia a zero dentro la banda derivata. È tutta
qui l'affermazione: non che qualche fase sia soddisfatta, ma che lo
sia la loro media.}
\label{fig:plugcycle}
\end{figure}
